\abstract % Структурный элемент: РЕФЕРАТ
В рамках работы разработан и апробирован программный комплекс для классификации Bitcoin-адресов по типу кошелька на основе датасета, сформированного в ходе предшествующих НИРС. Реализован полный пайплайн обработки данных: удаление высококоррелированных признаков, автоматический выбор стратегии устранения дисбаланса классов (SMOTE против ADASYN), стандартизация. Обучено шесть моделей машинного обучения: MLP, Random Forest, XGBoost, LightGBM, Stacking и kNN. Для деревьевидных моделей выполнена настройка гиперпараметров методом случайного поиска с кросс-валидацией. Проведена комплексная оценка качества: построены матрицы ошибок, ROC-кривые, вычислены взвешенная F1-мера и макро-усреднённый ROC-AUC. Для лучшей модели выполнена оптимизация порогов классификации по критерию Юдена и интерпретация предсказаний методом SHAP. Наилучший результат достигнут моделью XGBoost: взвешенная F1-мера~--- 0,7907; ROC-AUC~--- 0,9610.

Ключевые слова: Bitcoin, блокчейн, классификация адресов, машинное обучение, XGBoost, LightGBM, нейронная сеть, Stacking, SMOTE, SHAP.
